\section{Related Work}
\label{sec:related_work}

\subsection{Concept Prerequisite Learning}
Automated Concept Prerequisite Learning (CPL) aims to discover directed dependencies between educational concepts from text corpora \cite{chen2018prereq}. Early work by Liang et al.~\cite{liang2015prereq} introduced the Active Learning Concept Prerequisite Learning (AL-CPL) benchmark and proposed Reference Distance (RefD), measuring asymmetric cross-referencing between Wikipedia articles. While effective on hyperlinked encyclopedic networks, Wikipedia-based methods rely on dense citation graphs that are entirely absent in standard instructional textbooks.

Subsequent studies examined educational texts directly. Alzetta et al.~\cite{alzetta2020prereq,alzetta2021textbooks} investigated textbook prerequisite extraction across domains, observing that linear document order (temporal presentation prior) represents one of the strongest predictive signals in expository prose. Roy et al.~\cite{roy2019moocs} evaluated prerequisite extraction on video lecture transcripts and showed that unsupervised classifiers struggle with directional asymmetry unless explicit temporal constraints are enforced. Chen et al.~\cite{chen2018prereq} introduced topological graph constraints to enforce directed acyclic graph (DAG) properties on prerequisite trees. Our study builds upon these insights by directly evaluating whether modern open-weight language models provide measurable directional intelligence beyond the structural reading-order prior.

\subsection{LLMs for Knowledge Graph Extraction}
Recent research has increasingly applied Large Language Models to open-domain information extraction and knowledge graph construction \cite{carta2023iterative,trajanoska2023llm}. Commercial models such as GPT-3.5 and GPT-4 demonstrate competitive few-shot relation classification on standard NLP benchmarks (e.g., CoNLL04, TACRED). However, as Kommineni et al.~\cite{kommineni2024llm} observed, smaller open-weight models (sub-10B parameters) frequently suffer from schema non-conformance, directional hallucination, and prompt position sensitivity when deployed locally. 

Importantly, most existing evaluations of LLMs in knowledge extraction test encyclopedic relations (e.g., \textit{born-in}, \textit{subsidiary-of}) on general text. In educational domains, the core relation of interest—prerequisite dependency—is conceptual, asymmetric, and rarely signaled by explicit lexical markers. By conducting a controlled evaluation of \texttt{Qwen-2.5-7B} under swap-sampled prompting, our work isolates the actual directional reasoning capacity of locally deployable models on curriculum text.

\subsection{Structural Priors and Multi-Criteria Weighting}
Symbolic knowledge engineering has long utilized heuristic rules, such as Hearst patterns \cite{hearst1992hyponyms}, and multi-criteria decision models to assign relational confidence. A prominent formulation is Saaty's Analytic Hierarchy Process (AHP) \cite{saaty1980ahp}, which constructs pairwise comparison matrices across multiple textual features (e.g., lexical definition overlap, co-occurrence frequency, and section proximity). While the internal mathematical consistency of AHP matrices is well understood via eigenvalue decomposition ($CR < 0.10$), whether such multi-factor combinations yield empirical predictive power over simpler statistical counts has rarely been evaluated against external ground truth. Our study directly addresses this gap through query-ranking ablations on human-annotated prerequisite queries.
